% Part: lambda-calculus
% Chapter: introduction
% Section: beta

\documentclass[../../../include/open-logic-section]{subfiles}

\begin{document}

\olfileid{lam}{syn}{bet}
\olsection{$\beta$-হ্রাস}
% BN-SRC-289 TARGET-CORRECTION-BEGIN
\par\smallskip\noindent\textnormal{\small\textbf{উৎস-সংশোধন (BN-SRC-289):} ফাইলটি Syntax অধ্যায়ের ড্রাইভার থেকে আমদানি হয়, তার পথও \texttt{lambda-calculus/syntax/beta.tex}, এবং একই অধ্যায়ের অন্য পরিচয়গুলিও \texttt{syn} ব্যবহার করে; কিন্তু হিমায়িত রেন্ডার-পরিচয়টি \texttt{\string\olfileid\{lam\}\{int\}\{bet\}}। বাংলা লক্ষ্যে মধ্যবর্তী পরিচয়টি \texttt{syn} করা হয়েছে, যাতে আন্তঃউল্লেখটি প্রকৃত অধ্যায়-স্থানের সঙ্গে মেলে।} % BN-SRC-289 TARGET-CORRECTION-END

$(\lambd[m][(\lambd[y][y]) m])$ দেখলে স্বাভাবিকভাবেই অনুমান হয় যে
এর সঙ্গে $\lambd[m][m]$-এর কোনো সম্পর্ক আছে: দ্বিতীয় পদটি যেন প্রথম
পদটিকে “সরল” করার ফল। \emph{$\beta$-হ্রাস}-এর ধারণা এই স্বজ্ঞাটিকে
আনুষ্ঠানিকভাবে প্রকাশ করে।

\begin{defn}[$\beta$-সংকোচন, $\bredone$] \ollabel{defn:betacontr}
  \emph{$\beta$-সংকোচন} ($\bredone$) হলো পদগুলির উপর এমন ক্ষুদ্রতম
  সামঞ্জস্যশীল সম্বন্ধ যা নিচের শর্তটি পূরণ করে:
  \[
    (\lambd[x][N])Q \bredone \Subst{N}{Q}{x} 
  \]
  $P \bredone Q$ হলে বলি, $P$-কে \emph{$\beta$-সংকুচিত} করে $Q$
  পাওয়া যায়। $(\lambd[x][N])Q$ রূপের পদকে \emph{রিডেক্স} বলে।
\end{defn}

\begin{prob} \ollabel{prob:def}
  \olref[lam][syn][alp]{defn:aconvone}-এ বদ্ধ চলরাশির পরিবর্তনের জন্য
  যেমন করেছিলাম, তেমন করে $\beta$-সংকোচনের সমতুল্য আরোহমূলক সংজ্ঞাগুলি
  বিস্তারিতভাবে লেখো।
\end{prob}
  
\begin{defn}[$\beta$-হ্রাস, $\bred$] \ollabel{defn:betared}
  \emph{$\beta$-হ্রাস} ($\bred$) হলো পদগুলির উপর এমন ক্ষুদ্রতম
  প্রতিবিম্ব ধর্মবিশিষ্ট ও পরিযায়ী সম্বন্ধ যাতে~$\bredone$ অন্তর্ভুক্ত।
  $P \bred Q$ হলে বলি, $P$-কে $\beta$-হ্রাস করে $Q$ পাওয়া যায়।
\end{defn}

প্রসঙ্গ স্পষ্ট হলে $\bredone$-এর বদলে $\redone$, এবং $\bred$-এর বদলে
$\red$ লিখব।

অনানুষ্ঠানিকভাবে, $M \bred N$ ঠিক তখনই, যখন $\beta$-সংকোচনের শূন্য
বা একাধিক ধাপে $M$-কে $N$-এ বদলানো যায়।

\begin{defn}[$\beta$-স্বাভাবিক]
যে পদকে আর $\beta$-সংকুচিত করা যায় না, তাকে \emph{$\beta$-স্বাভাবিক}
বলে।
\end{defn}

$M \bred N$ এবং $N$ যদি $\beta$-স্বাভাবিক হয়, তবে $N$-কে~$M$-এর
একটি \emph{স্বাভাবিক রূপ} বলি। কোনো পদের স্বাভাবিক রূপ অনন্য কি না,
এই প্রশ্ন করা যায়; উত্তর হ্যাঁ, যা পরে দেখব।

কয়েকটি উদাহরণ দেখা যাক।
\begin{enumerate}
\item আমাদের আছে
\begin{align*}
(\lambd[x][xxy]) \lambd[z][z] & \redone (\lambd[z][z])(\lambd[z][z]) y \\
& \redone (\lambd[z][z]) y \\
& \redone y
\end{align*}
\item কোনো পদকে “সরল” করলে সেটি বাস্তবে আরও জটিলও হতে পারে:
\begin{align*}
(\lambd[x][xxy])(\lambd[x][xxy]) & \redone (\lambd[x][xxy])(\lambd[x][xxy])y \\
& \redone (\lambd[x][xxy])(\lambd[x][xxy])yy \\
& \redone \dots
\end{align*}
\item এতে কোনো পদ অপরিবর্তিতও থাকতে পারে:
\[
(\lambd[x][xx])(\lambd[x][xx]) \redone (\lambd[x][xx])(\lambd[x][xx])
\]
\item আবার কিছু পদকে একাধিকভাবে হ্রাস করা যায়। যেমন,
\[
(\lambd[x][(\lambd[y][yx]) z]) v \redone (\lambd[y][yv]) z
\]
সবচেয়ে বাইরের প্রয়োগটি সংকুচিত করলে এটি পাওয়া যায়; আর
\[
(\lambd[x][(\lambd[y][yx]) z]) v \redone (\lambd[x][zx]) v
\]
সবচেয়ে ভিতরেরটি সংকুচিত করলে এটি পাওয়া যায়। তবে লক্ষ করো, এই
ক্ষেত্রে উভয় পদই আরও হ্রাস পেয়ে একই পদ~$zv$ হয়।
\end{enumerate}

শেষ উদাহরণের চূড়ান্ত ফল কাকতালীয় নয়; বরং এটি ল্যাম্বডা ক্যালকুলাসের
একটি গভীর ও গুরুত্বপূর্ণ ধর্মের দৃষ্টান্ত, যা চার্চ--রোসার ধর্ম নামে
পরিচিত।

\begin{digress}
  সাধারণভাবে কোনো পদকে একাধিকভাবে $\beta$-হ্রাস করা যায়, তাই অনেক
  হ্রাস-কৌশল উদ্ভাবিত হয়েছে। এদের মধ্যে সবচেয়ে প্রচলিত হলো
  \emph{স্বাভাবিক কৌশল}। রিডেক্সের অবস্থান বলতে পদটির মধ্যে তার
  শুরুর স্থান বোঝালে, স্বাভাবিক কৌশল সর্বদা \emph{সর্ববাঁয়ের}
  রিডেক্সটিকে সংকুচিত করে। এই কৌশলের উপকারী ধর্ম হলো: কোনো কৌশলে
  একটি পদকে স্বাভাবিক রূপে হ্রাস করা যায় ঠিক তখনই, যখন স্বাভাবিক
  কৌশল ব্যবহার করেও তাকে স্বাভাবিক রূপে হ্রাস করা যায়। আলাদা করে
  কিছু না বললে পরবর্তী আলোচনায় আমরা স্বাভাবিক কৌশলই ব্যবহার করব।
\end{digress}

\begin{defn}[$\beta$-সমতুল্যতা, $\equal$]
  \emph{$\beta$-সমতুল্যতা} ($\equal$) হলো নিচের আরোহমূলকভাবে
  সংজ্ঞায়িত সম্বন্ধ:
  \begin{enumerate}
  \item $M \equal M$।
  \item $M \equal N$ হলে $N \equal M$।
  \item $M \equal N$, $N \equal O$ হলে $M \equal O$।
  \item $M \equal N$ হলে $PM \equal PN$।
  \item $M \equal N$ হলে $MQ \equal NQ$।
  \item $M \equal N$ হলে $\lambd[x][M] \equal \lambd[x][N]$।
  \item $(\lambd[x][N])Q \equal \Subst{N}{Q}{x}$।
  \end{enumerate}
\end{defn}

প্রথম তিনটি বিধি সম্বন্ধটিকে একটি তুল্যতা সম্পর্ক করে; পরের তিনটি
তাকে সামঞ্জস্যশীল করে; আর শেষ বিধিটি নিশ্চিত করে যে এতে
$\beta$-সংকোচন অন্তর্ভুক্ত।

অনানুষ্ঠানিকভাবে, দুটি পদ ঠিক তখনই $\beta$-সমতুল্য, যখন
$\beta$-সংকোচনের শূন্য বা একাধিক ধাপে, অথবা $\beta$-সংকোচনের
“বিপরীত” ধাপে, একটিকে অন্যটিতে বদলানো যায়। বিপরীত $\beta$-সংকোচন
এভাবে সংজ্ঞায়িত: $M$ থেকে বিপরীত-$\beta$-সংকোচনে $N$ পাওয়া যায়
ঠিক তখনই, যখন $N$ থেকে $\beta$-সংকোচনে $M$ পাওয়া যায়।

উপরের বিধিগুলির সঙ্গে পরে আরও বিধি যোগ করে সম্বন্ধটিকে প্রসারিত করব;
প্রসারিত তুল্যতা সম্পর্কটি $\equal[X]$ লিখব, যেখানে $X$ হলো
প্রসারণকারী বিধি।

\end{document}
